\section{Proofs for Section~\ref{sec:short_run}}

This appendix provides the algebra omitted from Section~\ref{sec:short_run}.

\subsection{Proof of Proposition \ref{prop:renewal_gap}}

Fix $h\in\{0,1\}$. From \eqref{eq:Pi_compact}, the decentralized benchmark solves
\[
\max_{b\ge 0}\left\{\bar\Pi(h)+\beta A_h^P b-\frac{\kappa}{2}b^2\right\}.
\]
Differentiating with respect to $b$ gives
\[
\frac{\partial \Pi(h,b)}{\partial b}=\beta A_h^P-\kappa b.
\]
Since
\[
\frac{\partial^2 \Pi(h,b)}{\partial b^2}=-\kappa<0,
\]
the problem is strictly concave, and the first-order condition is sufficient. Setting the derivative equal to zero yields
\[
b_P^*(h)=\frac{\beta A_h^P}{\kappa}
=
\frac{\beta v_XQ_{AB}(\eta_0+\eta_1 h)}{\kappa},
\]
which is \eqref{eq:bPstar}.

Likewise, from \eqref{eq:W_compact}, the planner solves
\[
\max_{b\ge 0}\left\{\bar W(h)+\beta A_h^W b-\frac{\kappa}{2}b^2\right\}.
\]
Differentiating,
\[
\frac{\partial W(h,b)}{\partial b}=\beta A_h^W-\kappa b,
\qquad
\frac{\partial^2 W(h,b)}{\partial b^2}=-\kappa<0,
\]
so the unique optimum is
\[
b_W^*(h)=\frac{\beta A_h^W}{\kappa}
=
\frac{\beta\Bigl[\xi v_LP_A m_A(h)+v_XQ_{AB}(\eta_0+\eta_1 h)\Bigr]}{\kappa},
\]
which is \eqref{eq:bWstar}.

Subtracting the two solutions gives
\[
b_W^*(h)-b_P^*(h)
=
\frac{\beta\bigl(A_h^W-A_h^P\bigr)}{\kappa}
=
\frac{\beta\xi v_LP_A m_A(h)}{\kappa}
>0,
\]
which proves \eqref{eq:renewal_gap}. Using $m_A(1)-m_A(0)=m_A^1-m_A^0>0$, we also obtain
\[
\bigl[b_W^*(1)-b_P^*(1)\bigr]-\bigl[b_W^*(0)-b_P^*(0)\bigr]
=
\frac{\beta\xi v_LP_A\bigl(m_A^1-m_A^0\bigr)}{\kappa}
>0,
\]
which is \eqref{eq:renewal_gap_hub}. This proves Proposition \ref{prop:renewal_gap}.
\qed

\subsection{Derivation of the Decentralized--Social Wedge}

Substituting the optimal pipeline choices into \eqref{eq:Pi_compact} and \eqref{eq:W_compact} gives
\[
\Pi^*(h)=\bar\Pi(h)+\frac{\beta^2(A_h^P)^2}{2\kappa}
=
\bar\Pi(h)+R_h^P,
\]
and
\[
W^*(h)=\bar W(h)+\frac{\beta^2(A_h^W)^2}{2\kappa}
=
\bar W(h)+R_h^W.
\]

Using \eqref{eq:Wbar} and \eqref{eq:Pi_compact},
\[
\bar W(1)-\bar W(0)
=
v_LP_A(m_A^1-m_A^0)
+
\beta v_LP_A\Bigl[m_A^1(1-\delta m_A^1)-m_A^0(1-\delta m_A^0)\Bigr]
-F.
\]
Factoring the second term,
\[
\bar W(1)-\bar W(0)=S_A+D_A-F,
\]
where $S_A$ and $D_A$ are defined in \eqref{eq:SA} and \eqref{eq:DA}. Also,
\[
\bar\Pi(1)-\bar\Pi(0)=S_A-F.
\]
Therefore
\[
\Delta_P^*=\Pi^*(1)-\Pi^*(0)=S_A-F+(R_1^P-R_0^P),
\]
which is \eqref{eq:DeltaP_star_expanded}, and
\[
\Delta_W^*=W^*(1)-W^*(0)=S_A+D_A-F+(R_1^W-R_0^W).
\]
Subtracting the two expressions gives
\[
\Delta_W^*-\Delta_P^*
=
D_A+\Bigl[(R_1^W-R_0^W)-(R_1^P-R_0^P)\Bigr]
=
D_A+M_A,
\]
which is \eqref{eq:DeltaW_star_wedge}.

To verify $M_A>0$, note that
\[
A_h^W-A_h^P=\xi v_LP_A m_A(h)>0
\qquad
\text{for } h\in\{0,1\},
\]
and that both $m_A(h)$ and $A_h^P$ are weakly larger when $h=1$ than when $h=0$. Hence
\[
(A_1^W)^2-(A_1^P)^2>(A_0^W)^2-(A_0^P)^2,
\]
which implies
\[
M_A
=
\frac{\beta^2}{2\kappa}
\left(
(A_1^W)^2-(A_1^P)^2-(A_0^W)^2+(A_0^P)^2
\right)
>0.
\]

\subsection{Proof of Proposition \ref{prop:dynamic_inefficiency}}

By definition,
\[
\Delta_P^*>0
\quad\Longleftrightarrow\quad
F<F_P^*,
\]
where
\[
F_P^*=S_A+(R_1^P-R_0^P).
\]
Likewise,
\[
\Delta_W^*>0
\quad\Longleftrightarrow\quad
F<F_W^*,
\]
where
\[
F_W^*=F_P^*+D_A+M_A.
\]

\noindent\textbf{Part 1.}
If $D_A+M_A>0$, then $F_W^*>F_P^*$. Therefore the interval
\[
F_P^*<F<F_W^*
\]
is nonempty. On this interval,
\[
\Delta_P^*<0
\qquad\text{and}\qquad
\Delta_W^*>0.
\]
Hence decentralized actors reject the hub while the planner supports it. This proves \eqref{eq:under_support_interval}.

\noindent\textbf{Part 2.}
If $D_A+M_A<0$, then $F_W^*<F_P^*$. Therefore the interval
\[
F_W^*<F<F_P^*
\]
is nonempty. On this interval,
\[
\Delta_W^*<0
\qquad\text{and}\qquad
\Delta_P^*>0.
\]
Hence decentralized actors support the hub while the planner rejects it. This proves \eqref{eq:dynamic_inefficiency_interval}.

This completes the proof of Proposition \ref{prop:dynamic_inefficiency}.
\qed
